\section*{Chapter \ref{ch:b1:complex} --- Complex Numbers}

\begin{solution}{exo:b1:complex:1}
$1 + \iu = \sqrt 2\, \eu^{\iu\pi/4}$;
$\;\sqrt 3 - \iu = 2\, \eu^{-\iu\pi/6}$;
$\;-5 = 5\, \eu^{\iu\pi}$;
\[
\frac{1 + \iu}{\sqrt 3 - \iu}
= \frac{\sqrt 2}{2}\, \eu^{\iu(\pi/4 + \pi/6)}
= \frac{\sqrt 2}{2}\, \eu^{5\iu\pi/12}.
\]
Computing the same quotient algebraically (multiply by the conjugate):
\[
\frac{(1 + \iu)(\sqrt 3 + \iu)}{4}
= \frac{(\sqrt 3 - 1) + \iu(\sqrt 3 + 1)}{4}.
\]
Identifying with $\frac{\sqrt 2}{2}(\cos\frac{5\pi}{12} +
\iu\sin\frac{5\pi}{12})$:
\[
\cos\frac{5\pi}{12} = \frac{\sqrt 6 - \sqrt 2}{4},
\qquad
\sin\frac{5\pi}{12} = \frac{\sqrt 6 + \sqrt 2}{4}.
\]
\end{solution}

\begin{solution}{exo:b1:complex:2}
$(1+\iu)^2 = 2\iu$, so $(1+\iu)^{20} = (2\iu)^{10} = 2^{10}\,\iu^{10}
= 1024 \times (-1) = -1024$.

In exponential form $(1+\iu)^n = 2^{n/2}\, \eu^{\iu n\pi/4}$, real if
and only if $\sin\frac{n\pi}{4} = 0$, i.e.\ $4 \mid n$. So $(1+\iu)^n
\in \R$ exactly for the multiples of $4$ (value $(-4)^{n/4}$).
\end{solution}

\begin{solution}{exo:b1:complex:3}
$\abs{z - 2} = \abs{z + \iu}$: equidistance from the points $2$ and
$-\iu$ --- the perpendicular bisector of the segment joining $(2, 0)$
and $(0, -1)$.

$\abs{z - 1} = 2$: the circle of center $1$ and radius $2$.

$\frac{z-1}{z+1} \in \iu\R$: by
\cref{prop:b1:complex:geometry}~(3), the points $M(z)$, $A(1)$,
$B(-1)$ satisfy: lines $MA$ and $MB$ perpendicular (or $z = 1$, where
the quotient is $0 \in \iu\R$). The set is the circle of diameter
$[-1, 1]$ (the unit circle), minus the point $-1$ where the quotient
is undefined. \emph{Check by computation:} $z = \eu^{\iu\theta}$ gives
$\frac{z-1}{z+1} = \frac{\eu^{\iu\theta/2}(\eu^{\iu\theta/2} -
\eu^{-\iu\theta/2})}{\eu^{\iu\theta/2}(\eu^{\iu\theta/2} +
\eu^{-\iu\theta/2})} = \iu\tan\frac\theta2 \in \iu\R$.
\end{solution}

\begin{solution}{exo:b1:complex:4}
Linearization:
\[
\sin^4\theta
= \Bigl(\frac{\eu^{\iu\theta} - \eu^{-\iu\theta}}{2\iu}\Bigr)^{\!4}
= \frac{\eu^{4\iu\theta} - 4\eu^{2\iu\theta} + 6 - 4\eu^{-2\iu\theta}
+ \eu^{-4\iu\theta}}{16}
= \frac{\cos 4\theta - 4\cos 2\theta + 3}{8}.
\]
Expansion: by de Moivre, $\cos 4\theta = \Re\bigl((c +
\iu s)^4\bigr) = c^4 - 6c^2 s^2 + s^4$ with $c = \cos\theta$, $s =
\sin\theta$; substituting $s^2 = 1 - c^2$:
\[
\cos 4\theta = c^4 - 6c^2(1 - c^2) + (1 - c^2)^2
= 8c^4 - 8c^2 + 1 .
\]
\end{solution}

\begin{solution}{exo:b1:complex:5}
$8\iu = 8\,\eu^{\iu\pi/2}$, so the cube roots are $2\,\eu^{\iu(\pi/6 +
2k\pi/3)}$, $k = 0, 1, 2$:
\[
z_0 = 2\eu^{\iu\pi/6} = \sqrt 3 + \iu,\quad
z_1 = 2\eu^{5\iu\pi/6} = -\sqrt 3 + \iu,\quad
z_2 = 2\eu^{3\iu\pi/2} = -2\iu :
\]
an equilateral triangle on the circle of radius $2$.

$-4 = 4\eu^{\iu\pi}$, so $z^4 = -4$ has solutions $\sqrt 2\,
\eu^{\iu(\pi/4 + k\pi/2)}$: $\;1 + \iu$, $-1 + \iu$, $-1 - \iu$,
$1 - \iu$. Pairing conjugate roots:
\[
X^4 + 4 = \bigl(X^2 - 2X + 2\bigr)\bigl(X^2 + 2X + 2\bigr),
\]
since $(X - (1+\iu))(X - (1-\iu)) = X^2 - 2X + 2$ and similarly for
the other pair. (Expand to check.)
\end{solution}

\begin{solution}{exo:b1:complex:6}
$C_n + \iu S_n = \sum_{k=0}^{n} \eu^{\iu k\theta} =
\frac{\eu^{\iu(n+1)\theta} - 1}{\eu^{\iu\theta} - 1}$ (geometric sum,
ratio $\eu^{\iu\theta} \neq 1$). Half-angle factorization
(\cref{met:b1:complex:trig}~(4)) on numerator and denominator:
\[
\frac{2\iu\sin\frac{(n+1)\theta}{2}\, \eu^{\iu(n+1)\theta/2}}
{2\iu\sin\frac{\theta}{2}\, \eu^{\iu\theta/2}}
= \frac{\sin\frac{(n+1)\theta}{2}}{\sin\frac{\theta}{2}}\,
\eu^{\iu n\theta/2}.
\]
Taking real and imaginary parts:
\[
C_n = \frac{\sin\frac{(n+1)\theta}{2}}{\sin\frac{\theta}{2}}
\cos\frac{n\theta}{2},
\qquad
S_n = \frac{\sin\frac{(n+1)\theta}{2}}{\sin\frac{\theta}{2}}
\sin\frac{n\theta}{2}.
\]
\end{solution}

\begin{solution}{exo:b1:complex:7}
Discriminant: $\Delta = (3 + 4\iu)^2 - 4(-1 + 5\iu) = 9 + 24\iu - 16 +
4 - 20\iu = -3 + 4\iu$. Square roots of $-3 + 4\iu$: solve $x^2 - y^2
= -3$, $2xy = 4$, $x^2 + y^2 = 5$; then $x^2 = 1$, $y = 2/x$, same
signs: $\delta = \pm(1 + 2\iu)$. Hence
\[
z = \frac{(3 + 4\iu) \pm (1 + 2\iu)}{2}
\in \{\, 2 + 3\iu,\; 1 + \iu \,\}.
\]
\emph{Check:} sum $= 3 + 4\iu$ and product $(2+3\iu)(1+\iu) = -1 +
5\iu$, as required by the coefficients.
\end{solution}

\begin{solution}{exo:b1:complex:8}
\begin{enumerate}
  \item \cref{prop:b1:complex:sumroots} with $n = 5$.
  \item $u + v = \omega + \omega^2 + \omega^3 + \omega^4 = -1$ by (1).
        For the product, expand and reduce exponents modulo $5$:
        \[
        uv = (\omega + \omega^4)(\omega^2 + \omega^3)
        = \omega^3 + \omega^4 + \omega^6 + \omega^7
        = \omega^3 + \omega^4 + \omega + \omega^2 = -1 .
        \]
        So $u$ and $v$ have sum $-1$ and product $-1$: they are the
        two roots of $X^2 + X - 1$.
  \item $u = \omega + \conj\omega = 2\cos\frac{2\pi}{5} > 0$ (the
        angle is acute), and the positive root of $X^2 + X - 1$ is
        $\frac{-1 + \sqrt 5}{2}$. Hence $\cos\frac{2\pi}{5} =
        \frac{\sqrt 5 - 1}{4}$.
\end{enumerate}
\end{solution}

\begin{solution}{exo:b1:complex:9}
Expand both squares as in the proof of
\cref{prop:b1:complex:rules}~(4):
\[
\abs{z + w}^2 = \abs z^2 + \abs w^2 + 2\Re(z\conj w),
\qquad
\abs{z - w}^2 = \abs z^2 + \abs w^2 - 2\Re(z\conj w),
\]
and add. Geometrically, $\abs{z+w}$ and $\abs{z-w}$ are the lengths of
the two diagonals of the parallelogram with vertices $0, z, z+w, w$,
while $\abs z$ and $\abs w$ are the side lengths: the sum of the
squares of the diagonals equals the sum of the squares of the four
sides.
\end{solution}

\begin{solution}{exo:b1:complex:10}
Let $j = \eu^{2\iu\pi/3}$, so $j^2 + j + 1 = 0$ and $\eu^{\iu\pi/3} =
-j^2$, $\eu^{-\iu\pi/3} = -j$. The triangle is equilateral direct iff
$c - a = -j^2 (b - a)$, indirect iff $c - a = -j(b - a)$
(\cref{ex:b1:complex:equilateral}).

Consider $P = a + jb + j^2 c$ and $Q = a + j^2 b + jc$. Using $1 +
j^2 = -j$ and $j^3 = 1$:
\[
c - a + j^2(b - a) = c + j^2 b - (1 + j^2)\,a = c + j^2 b + ja
= j\,(a + jb + j^2 c) = jP,
\]
so the direct case reads $jP = 0 \iff P = 0$; the same computation
with $j$ in place of $j^2$ gives $c - a + j(b - a) = j^2 Q$, so the
indirect case reads $Q = 0$. Hence: equilateral (either orientation)
$\iff PQ = 0$. Expanding, with $j + j^2 = -1$:
\[
PQ = a^2 + b^2 + c^2 + (j + j^2)(ab + bc + ca)
= a^2 + b^2 + c^2 - (ab + bc + ca).
\]
So the triangle is equilateral if and only if $a^2 + b^2 + c^2 = ab +
bc + ca$.
\end{solution}

\begin{solution}{exo:b1:complex:11}
Since the $\omega^k$, $k = 0, \dots, n-1$, are exactly the $n$ roots
of $X^n - 1$ (\cref{thm:b1:complex:roots}), and the polynomial is
monic:
\[
X^n - 1 = \prod_{k=0}^{n-1} (X - \omega^k)
= (X - 1) \prod_{k=1}^{n-1} (X - \omega^k).
\]
Dividing by $X - 1$: $\;1 + X + \dots + X^{n-1} = \prod_{k=1}^{n-1}
(X - \omega^k)$. Evaluating at $X = 1$ gives $P = n$.

Now $1 - \omega^k = -\eu^{\iu k\pi/n}\bigl(\eu^{\iu k\pi/n} -
\eu^{-\iu k\pi/n}\bigr) = -2\iu\,\eu^{\iu k\pi/n} \sin\frac{k\pi}{n}$,
so taking moduli in $P = n$ (each $\sin\frac{k\pi}n > 0$ for $1 \leq k
\leq n-1$):
\[
n = \abs P = \prod_{k=1}^{n-1} 2\sin\frac{k\pi}{n}
= 2^{n-1} \prod_{k=1}^{n-1} \sin\frac{k\pi}{n},
\qquad\text{hence}\qquad
\prod_{k=1}^{n-1} \sin\frac{k\pi}{n} = \frac{n}{2^{n-1}} .
\]
\end{solution}
